\chapter{Cost Functions}
The `costMat` function serves as our primary heuristic, allowing for the selection of five distinct cost metrics. These metrics—logarithmic, linear sum, square, cubic, and quartic—are controlled by a parameter $p$, which determines the mathematical weight assigned to the matrix reduction process.

\section{Log}
This cost function is derived from the greedy framework proposed by \cite{9475957}, which utilizes the logarithm of the Hamming weight for each row as follows:
$$h_{log}(A) = \sum_{i}log_{2}(\sum_{j}a_{i,j})$$

\section{Sum}
In addition to the logarithmic metric, the greedy framework by \cite{9475957} considers the cumulative Hamming weight of each row as follows:
$$h_{sum}(A) = \sum_{i}\sum_{j}a_{i,j}$$

\section{Square}
Following the methodology of \cite{cryptoeprint:2024/381}, this cost function utilizes the square of each row's Hamming weight, defined as follows:
$$h_{sq}(A) = \sum_{i}(\sum_{j}a_{i,j})^{2}$$

\section{The Role of Matrix Cost Functions in Greedy Optimization}
The `costMat` function is integrated into two distinct phases of our greedy algorithm: the initial candidate selection and the subsequent layer evaluation.

\subsection{Globally}
First, the algorithm evaluates the global maximum cost to establish a relaxation range; this allows the search to include a broader set of available operators rather than being restricted to the absolute minimum.

\begin{algorithm}
\caption{Greedy}
\begin{algorithmic}[1]
\State $\vdots$
\While{$mat$ != Permutation}
	\State $\vdots$
	\If{p != 1}
	\State $H_r$ $\leftarrow$ costMat($mat$, $p$) + costMat($inv^{T}$, $p$)
	\State $H_c$ $\leftarrow$ costMat($mat^{T}$, $p$) + costMat($inv$, $p$)
	\State minm $\leftarrow$ max\{$H_r$, $H_c$\}
\Else
	\State minm $\leftarrow$ costMat($mat$, $p$) + costMat($inv$, $p$)
\EndIf
\State $\vdots$
\EndWhile
\end{algorithmic}
\end{algorithm}

Following the greedy framework established by Shi and Feng, we define the cost metrics $H_{r}$ and $H_{c}$ to treat row and column operations asymmetrically. Specifically, the row Hamming weight is evaluated when performing row operations ($row_{op}$), while the column Hamming weight is prioritized during column operations ($col_{op}$). 

To explore the set of linear reversible operators for synthesizing quantum cryptographic circuits, we employ the following guidance:
$$H_{r} = costMat(A, p) + costMat((A^{-1})^{T}, p)$$
$$H_{c} = costMat(A^{T}, p) + costMat(A^{-1}, p)$$

\subsection{Locally}
In the second phase, the algorithm performs a local evaluation to identify individual operators that minimize the cost function. If a candidate operator's cost is lower than the current minimum, it is accepted as a valid transformation. 

Once all potential control and target qubit pairs have been visited, the resulting set of parallel operations is appended to the `select\_list` as a single circuit layer.

The availableRowOperatorSelection and availableColOperatorSelection functions influence the cost matrix calculation in the Row and Col greedy algorithms.

\begin{algorithm}
\caption{availableRowOperatorSelection}
\begin{algorithmic}[1]
\Require $L_{row}, L_{row}^{cost}, mat_{n \times n}$, $inv_{n \times n}, p, minm$ and select\_list
\Ensure select\_list, $minm$
\For{$row_{op} \in L_{row}$}
	\State temp\_mat $\leftarrow$ row\_i2j($mat$, $row_{op}$[0], $row_{op}$[1])
	\State temp\_inv $\leftarrow$ col\_i2j($inv$, $row_{op}$[1], $row_{op}$[0])
	\If{p != 1}
		\State $L_{row}^{cost}$.append(costMat(temp\_mat, p)+costMat(temp\_inv$^{T}$, p))
	\Else
		\State $L_{row}^{cost}$.append(costMat(temp\_mat, p)+costMat(temp\_inv, p))
	\EndIf
\EndFor
\For{index, $row_{op}^{cost}$ in $L_{row}^{cost}$}
	\If{$row_{op}^{cost} < minm$}
		\State select\_list.clear()
		\State select\_list.append(($L_{row}$[index][0], $L_{row}$[index][1], 0))
		\State $minm = row_{op}^{cost}$
	\ElsIf{$row_{op}^{cost} == minm$}
		\State select\_list.append(($L_{row}$[index][0], $L_{row}$[index][1], 0))
    \EndIf
\EndFor
\State \Return select\_list, $minm$
\end{algorithmic}
\end{algorithm}

\begin{algorithm}
\caption{availableColOperatorSelection}
\begin{algorithmic}[1]
\Require $L_{col}, L_{col}^{cost}, mat_{n \times n}$, $inv_{n \times n}, p, minm$ and select\_list
\Ensure select\_list, $minm$
\For{$col_{op} \in L_{col}$}
	\State temp\_mat $\leftarrow$ col\_i2j($mat$, $col_{op}$[0], $col_{op}$[1])
	\State temp\_inv $\leftarrow$ row\_i2j($inv$, $col_{op}$[1], $col_{op}$[0])
	\If{p != 1}
		\State $L_{col}^{cost}$.append(costMat(temp\_mat$^{T}$, p)+costMat(temp\_inv, p))
	\Else
		\State $L_{col}^{cost}$.append(costMat(temp\_mat, p)+costMat(temp\_inv, p))
	\EndIf
\EndFor
\For{index, $col_{op}^{cost}$ in $L_{col}^{cost}$}
	\If{$col_{op}^{cost} < minm$}
		\State select\_list.clear()
		\State select\_list.append(($L_{col}$[index][0], $L_{col}$[index][1], 1))
		\State $minm = col_{op}^{cost}$
	\ElsIf{$col_{op}^{cost} == minm$}
		\State select\_list.append(($L_{col}$[index][0], $L_{col}$[index][1], 1))
    \EndIf
\EndFor
\State \Return select\_list, $minm$
\end{algorithmic}
\end{algorithm}

Similarly, the modified and avoidLocalMinima versions of these functions influence the cost matrix in the Row-or-Col and Parallel greedy algorithms.

For the Row-or-Column and Parallel strategies, which evaluate two directional operators simultaneously, the algorithm selects the layer that minimizes the total matrix cost. 

To facilitate this, we implement an intermediate storage list for candidate operations and utilize their respective cost impacts as the primary selection criterion. 

This represents a departure from standard operator selection, as detailed in the following modifications:

\begin{algorithm}
\caption{modifiedAvailableRowOperatorSelection}
\begin{algorithmic}[1]
\Require $L_{row}, L_{row}^{cost}, mat_{n \times n}$, $inv_{n \times n}, p, minm, best_{row}^{cost}$ and $B_{row}$
\Ensure $B_{row}, best_{row}^{cost}, minm$
\For{$row_{op} \in L_{row}$}
	\State temp\_mat $\leftarrow$ row\_i2j($mat$, $row_{op}$[0], $row_{op}$[1])
	\State temp\_inv $\leftarrow$ col\_i2j($inv$, $row_{op}$[1], $row_{op}$[0])
	\If{p != 1}
		\State $L_{row}^{cost}$.append(costMat(temp\_mat, p)+costMat(temp\_inv$^{T}$, p))
	\Else
		\State $L_{row}^{cost}$.append(costMat(temp\_mat, p)+costMat(temp\_inv, p))
	\EndIf
\EndFor
\For{index, $row_{op}^{cost}$ in $L_{row}^{cost}$}
	\If{$row_{op}^{cost} < minm$}
		\If{$row_{op}^{cost} < best_{row}^{cost}$}
			\State $B_{row}$.clear()
			\State $best_{row}^{cost} = row_{op}^{cost}$
			\If{$row_{op}^{cost} == best_{row}^{cost}$}
				\State $B_{row}$.append(($L_{row}$[index][0], $L_{row}$[index][1], 0))
            \EndIf
        \EndIf
    \EndIf
\EndFor
\State \Return $B_{row}, best_{row}^{cost}$
\end{algorithmic}
\end{algorithm}

\begin{algorithm}
\caption{modifiedAvailableColOperatorSelection}
\begin{algorithmic}[1]
\Require $L_{col}, L_{col}^{cost}, mat_{n \times n}$, $inv_{n \times n}, p, minm, best_{col}^{cost}$ and $B_{col}$
\Ensure $B_{col}, best_{col}^{cost}, minm$
\For{$col_{op} \in L_{col}$}
	\State temp\_mat $\leftarrow$ col\_i2j($mat$, $col_{op}$[0], $col_{op}$[1])
	\State temp\_inv $\leftarrow$ row\_i2j($inv$, $col_{op}$[1], $col_{op}$[0])
	\If{p != 1}
		\State $L_{col}^{cost}$.append(costMat(temp\_mat$^{T}$, p)+costMat(temp\_inv, p))
	\Else
		\State $L_{col}^{cost}$.append(costMat(temp\_mat, p)+costMat(temp\_inv, p))
	\EndIf
\EndFor
\For{index, $col_{op}^{cost}$ in $L_{col}^{cost}$}
	\If{$col_{op}^{cost} < minm$}
		\If{$col_{op}^{cost} < best_{col}^{cost}$}
			\State $B_{col}$.clear()
			\State $best_{col}^{cost} = col_{op}^{cost}$
			\If{$col_{op}^{cost} == best_{col}^{cost}$}
				\State $B_{col}$.append(($L_{col}$[index][0], $L_{col}$[index][1], 1))
            \EndIf
        \EndIf
    \EndIf
\EndFor
\State \Return $B_{col}, best_{col}^{cost}$
\end{algorithmic}
\end{algorithm}

To address the issue of local minima, we implemented a stagnation recovery mechanism. When the algorithm becomes 'stuck'—meaning no cost-reducing operations are found—the greedy constraint is temporarily relaxed. 

During these instances, the algorithm accepts operations that slightly increase the current minimum cost. 

This heuristic perturbation enables the search to bypass stagnant states and discover subsequent operations that lead to a more optimal global reduction.

\begin{algorithm}
\caption{avoidLocalMinimaRowOperatorSelection}
\begin{algorithmic}[1]
\Require $L_{row}, mat_{n \times n}$, $inv_{n \times n}, p$ and escapeCandidates
\Ensure escapeCandidates
\For{$row_{op} \in L_{row}$}
	\State temp\_mat $\leftarrow$ row\_i2j($mat$, $row_{op}$[0], $row_{op}$[1])
	\State temp\_inv $\leftarrow$ col\_i2j($inv$, $row_{op}$[1], $row_{op}$[0])
	\If{p != 1}
		\State temp\_cost $\leftarrow$ costMat(temp\_mat, p)+costMat(temp\_inv$^{T}$, p)
	\Else
		\State temp\_cost $\leftarrow$ costMat(temp\_mat, p)+costMat(temp\_inv, p)
	\EndIf
	\State escapeCandidates.append((temp\_cost, $row_{op}$[0], $row_{op}$[1], 0))
\EndFor
\State \Return escapeCandidates
\end{algorithmic}
\end{algorithm}

\begin{algorithm}
\caption{avoidLocalMinimaColOperatorSelection}
\begin{algorithmic}[1]
\Require $L_{col}, mat_{n \times n}$, $inv_{n \times n}, p$ and escapeCandidates
\Ensure escapeCandidates
\For{$col_{op} \in L_{col}$}
	\State temp\_mat $\leftarrow$ col\_i2j($mat$, $col_{op}$[0], $col_{op}$[1])
	\State temp\_inv $\leftarrow$ row\_i2j($inv$, $col_{op}$[1], $col_{op}$[0])
	\If{p != 1}
		\State temp\_cost $\leftarrow$ costMat(temp\_mat$^{T}$, p)+costMat(temp\_inv, p)
	\Else
		\State temp\_cost $\leftarrow$ costMat(temp\_mat, p)+costMat(temp\_inv, p)
	\EndIf
	\State escapeCandidates.append((temp\_cost, $col_{op}$[0], $col_{op}$[1], 1))
\EndFor
\State \Return escapeCandidates
\end{algorithmic}
\end{algorithm}